
% VERSION 3
\subsection{Machine Unlearning Methods}
\label{subsec:gradient_unlearning}

Gradient Ascent (GA;~\citealt{jang2022knowledge}) maximizes the negative log-likelihood on the forget set without constraint, frequently causing catastrophic collapse of linguistic coherence~\cite{yao2024large, zhang2024negative}. Gradient Difference (GD;~\citealt{liu2022continual}) bounds collateral damage with a retain loss but assigns equal gradient weight to every forget token~\cite{bu2025unlearning}. Negative Preference Optimization (NPO;~\citealt{zhang2024negative}) frames unlearning as preference alignment, pushing the output distribution away from the memorized response. Weighted Gradient Ascent (WGA;~\citealt{wang2025rethinking}) weights each token's gradient by the model's log-probability, concentrating unlearning on high-confidence tokens. All four calibrate the unlearning signal from the model's current output distribution, which changes throughout training and does not reflect the degree to which a fact is parametrically encoded. \citet{lucki2025adversarial} further show that inference-time interventions leave parametric knowledge intact, with models remaining vulnerable to adversarial extraction under targeted queries.

\subsection{Memorization, Frequency, and Unlearning}
\label{subsec:memorization}

The extent of LLM memorization scales with training-corpus frequency~\cite{carlini2023quantifying, tirumala2022memorization}: \citet{kandpal2023large} show LLMs recall long-tail facts far less reliably than frequent ones, and \citet{merullo2025linear} demonstrate that frequent facts are encoded more linearly in model representations, providing a mechanistic basis for why popular facts resist gradient-based removal. This asymmetry is adversarial to unlearning: verbatim memorization impedes removal in both vision~\cite{NEURIPS2024_16e18fa3} and LLM settings~\cite{barbulescu2024each}, and \citet{baluta2024unlearning} confirm in a controlled setup that differently exposed data points are unlearned differently. Log-probabilities do not accurately reflect memorization depth~\cite{wang2025rethinking, hartmann2024undesirable}: a model may assign high confidence to a surface pattern even when the corresponding fact is not deeply encoded. External popularity proxies, such as entity citelink counts, provide a stable fact-level signal independent of the model's output distribution.

\subsection{Popularity Effects on Unlearning}
\label{subsec:popularity_effects}

\citet{krishnannot} establish through controlled pretraining experiments that high-frequency facts are either not unlearned or only superficially forgotten, remaining extractable via adversarial queries across multiple methods and model scales. \citet{anna2026anatomy} confirm the same failure mode on real-world facts in the DUET benchmark, showing that Wikipedia-based fact salience systematically predicts unlearning difficulty across SOTA unlearning methods. Both works precisely characterize the popularity gap and provide benchmarks we build on, but neither proposes a training method to address it. AdaPop is, to our knowledge, the first unlearning method that conditions per-fact gradient updates on an external popularity signal.





%% PLAN FROM 14/03/26
% \subsection{Gradient-Based Machine Unlearning}
% Weight-editing methods erase target facts by directly modifying model parameters, motivated by privacy regulations and the right to erasure~\cite{yao2024large}. Early approaches relied on prompting-based or in-context interventions~\cite{liu2024survey}; \citet{lucki2025adversarial} show that such inference-time methods leave parametric knowledge intact, with models remaining vulnerable to adversarial extraction under targeted queries. Parameter-editing methods address this gap but introduce a different challenge: none accounts for variation in how deeply individual facts are encoded across the model.

% Gradient Ascent (GA) maximizes the negative log-likelihood on the forget set, which frequently leads to catastrophic collapse of linguistic coherence~\cite{yao2024large, zhang2024negative}. Gradient Difference (GD) adds a retain loss to bound collateral damage but treats every token in the forget set with equal importance~\cite{bu2025unlearning}. Negative Preference Optimization (NPO) frames unlearning as preference alignment, pushing the output distribution away from the memorised response~\cite{zhang2024negative}. Weighted Gradient Ascent (WGA) weights each token's gradient by the model's log-probability, concentrating unlearning effort on high-confidence tokens~\cite{wang2025rethinking}. All four methods calibrate the unlearning signal from the model's current output distribution, which changes throughout training and does not reflect how deeply a fact is parametrically encoded~\cite{zhang2024negative, wang2025rethinking}. AdaPop replaces this local signal with a global popularity proxy, directly addressing the per-fact calibration gap that none of these methods resolves.

% \subsection{Knowledge Popularity and Memorization}
% \citet{kandpal2023large} show that LLMs recall long-tail facts far less reliably than frequent ones, and \citet{hartmann2024undesirable} survey how parametric memorisation scales with training-corpus frequency. Popular facts are redundantly encoded across many layers; rare facts are lightly encoded and susceptible to over-erasure when subjected to the same gradient. This asymmetry means a single gradient weight cannot address both regimes simultaneously.

% наиболее близкие к нам работы в MU - {\cite{krishnannot, anna2026anatomy}}. Они оценили сота методы такие как Х1 Х2 Х3 и экспериментально подвердили фандинг проблем. Однако, не предлагали свое решение. To the best of our knowledge, our proposed method is the first pop-aware MU method to fix the issue. 

% Log-probabilities do not accurately reflect memorisation depth~\cite{zhang2024negative}: a model may assign high confidence to a surface token pattern even when the corresponding fact is not deeply encoded across layers. External popularity proxies, such as entity citelink counts or co-occurrence frequencies, provide a fact-level signal independent of the model's current output distribution {\cite{krishnannot, anna2026anatomy}}. 

% %AdaPop assigns each fact a popularity-dependent exponent, treating memorisation depth as a property of the fact rather than of the model.
